% GENERATED FILE. Edit canonical lessons, then run npm run generate:books.
% canonical-source-hash: fnv1a64:21b59a78831b6c05
% canonical-lessons: JA-C515-juken-suru, JA-C515-ukaru, JA-C515-goukaku-suru, JA-C515-seiseki, JA-C515-yobikou, JA-R515-review-510-515

\chapter{If I Pass}
\label{ch:ja-if-i-pass}
\hlchaptermodality{515}

\begin{chapteropening}
\textbf{By the end of this chapter:} \emph{I can say what will happen if I pass or fail an exam.}

Complete the last lesson of chapter 515: I can say what will happen if I pass or fail an exam.
\end{chapteropening}

\section[\texorpdfstring{\ja{じゅけんする} (juken suru)}{じゅけんする (juken suru)}]{\ja{じゅけんする} (juken suru) — to sit an entrance exam}

\label{lesson:JA-C515-juken-suru}



\begin{quote}
Before the new one: say the Japanese for a forecast, then the Japanese for an outlook.
\end{quote}

\subsection*{You'll want to know: \ja{じゅけんする}}
\textbf{\ja{じゅけんする}} — \emph{juken suru} — \textquotedblleft{}to sit an entrance exam\textquotedblright{}.

\begin{quote}
\textbf{\ja{もし} \ja{らいねん} \ja{じゅけんしたら、まいにち} \ja{べんきょうします。}} — \emph{moshi rainen jukenshitara, mainichi benkyō shimasu} — \textquotedblleft{}If I sit the entrance exam next year, I will study every day.\textquotedblright{}
\end{quote}

Polite \textbf{\ja{じゅけんします}}, \emph{jukenshimasu}; a \ja{する} verb. Past (polite): \textbf{\ja{むすこが} \ja{ふたつの} \ja{がっこうを} \ja{じゅけんしました。}}, \emph{musuko ga futatsu no gakkō o jukenshimashita}.

\begin{cousinweb}
\textbf{\ja{じゅけんする}} (\emph{juken suru}) is the Sino-Japanese (kango) noun \emph{juken} with the native verb \emph{suru}, \textquotedblleft{}to do\textquotedblright{}: two Chinese characters read with their on'yomi, readings borrowed from Chinese. They are \emph{ju}, \textquotedblleft{}accept\textquotedblright{} (Middle Chinese \emph{dzyuwX}) and \emph{ken}, \textquotedblleft{}verification\textquotedblright{} (\emph{ngjemH}) (KANJIDIC2; Middle Chinese in Baxter's transcription).
\end{cousinweb}

\begin{grammarlens}[title={What you've built}]
One more word for saying what happens if.
\end{grammarlens}

\subsection*{Your turn}
\begin{itemize}
\raggedright
  \item \emph{Say it:} \emph{juken suru}
  \item \emph{Say it:} \emph{juken suru}, once more
  \item \emph{Say it:} \emph{moshi rainen jukenshitara, mainichi benkyō shimasu} — If I sit the entrance exam next year, I will study every day.
  \item \emph{Recall:} say the Japanese for a forecast, then the Japanese for an outlook, then say \emph{juken suru} again
\end{itemize}

\begin{tcolorbox}[breakable,colback=teal!4,colframe=teal!35!black,title={Before you move on}]
What does \ja{じゅけんする} mean? (\textquotedblleft{}To sit an entrance exam\textquotedblright{}.) Say it once more.
\end{tcolorbox}
\section[\texorpdfstring{\ja{うかる} (ukaru)}{うかる (ukaru)}]{\ja{うかる} (ukaru) — to pass (an exam)}

\label{lesson:JA-C515-ukaru}



\begin{quote}
Before the new one: say the Japanese for an outlook, then the Japanese for to sit an entrance exam.
\end{quote}

\subsection*{You'll want to know: \ja{うかる}}
\textbf{\ja{うかる}} — \emph{ukaru} — \textquotedblleft{}to pass (an exam)\textquotedblright{}.

\begin{quote}
\textbf{\ja{もし} \ja{しけんに} \ja{うかったら、みんなで} \ja{いわいます。}} — \emph{moshi shiken ni ukattara, minna de iwaimasu} — \textquotedblleft{}If I pass the exam, we will all celebrate.\textquotedblright{}
\end{quote}

Polite \textbf{\ja{うかります}}, \emph{ukarimasu}; a godan verb. Past (polite): \textbf{\ja{あねが} \ja{しけんに} \ja{うかりました。}}, \emph{ane ga shiken ni ukarimashita}.

\begin{cousinweb}
\textbf{\ja{うかる}} (\emph{ukaru}) is a native Japanese word (wago), inherited rather than borrowed from Chinese; its kanji was attached later as a kun'yomi, and its older origin is not settled (JMdict, KANJIDIC2).
\end{cousinweb}

\begin{grammarlens}[title={What you've built}]
One more word for saying what happens if.
\end{grammarlens}

\subsection*{Your turn}
\begin{itemize}
\raggedright
  \item \emph{Say it:} \emph{ukaru}
  \item \emph{Say it:} \emph{ukaru}, once more
  \item \emph{Say it:} \emph{moshi shiken ni ukattara, minna de iwaimasu} — If I pass the exam, we will all celebrate.
  \item \emph{Recall:} say the Japanese for an outlook, then the Japanese for to sit an entrance exam, then say \emph{ukaru} again
\end{itemize}

\begin{tcolorbox}[breakable,colback=teal!4,colframe=teal!35!black,title={Before you move on}]
What does \ja{うかる} mean? (\textquotedblleft{}To pass (an exam)\textquotedblright{}.) Say it once more.
\end{tcolorbox}
\section[\texorpdfstring{\ja{ごうかくする} (gōkaku suru)}{ごうかくする (gōkaku suru)}]{\ja{ごうかくする} (gōkaku suru) — to pass (a test); to be accepted}

\label{lesson:JA-C515-goukaku-suru}



\begin{quote}
Before the new one: say the Japanese for to sit an entrance exam, then the Japanese for to pass (an exam).
\end{quote}

\subsection*{You'll want to know: \ja{ごうかくする}}
\textbf{\ja{ごうかくする}} — \emph{gōkaku suru} — \textquotedblleft{}to pass (a test); to be accepted\textquotedblright{}.

\begin{quote}
\textbf{\ja{もし} \ja{ごうかくしたら、りょうしんは} \ja{よろこぶでしょう。}} — \emph{moshi gōkakushitara, ryōshin wa yorokobu deshō} — \textquotedblleft{}If I pass, my parents will surely be glad.\textquotedblright{}
\end{quote}

Polite \textbf{\ja{ごうかくします}}, \emph{gōkakushimasu}; a \ja{する} verb. Past (polite): \textbf{\ja{むすめが} \ja{ごうかくしました。}}, \emph{musume ga gōkakushimashita}.

\begin{cousinweb}
\textbf{\ja{ごうかくする}} (\emph{gōkaku suru}) is the Sino-Japanese (kango) noun \emph{gōkaku} with the native verb \emph{suru}, \textquotedblleft{}to do\textquotedblright{}: two Chinese characters read with their on'yomi, readings borrowed from Chinese. They are \emph{gō}, \textquotedblleft{}fit\textquotedblright{} (Middle Chinese \emph{hop}) and \emph{kaku}, \textquotedblleft{}status\textquotedblright{} (\emph{kaek}) (KANJIDIC2; Middle Chinese in Baxter's transcription).
\end{cousinweb}

\begin{grammarlens}[title={What you've built}]
One more word for saying what happens if.
\end{grammarlens}

\subsection*{Your turn}
\begin{itemize}
\raggedright
  \item \emph{Say it:} \emph{gōkaku suru}
  \item \emph{Say it:} \emph{gōkaku suru}, once more
  \item \emph{Say it:} \emph{moshi gōkakushitara, ryōshin wa yorokobu deshō} — If I pass, my parents will surely be glad.
  \item \emph{Recall:} say the Japanese for to sit an entrance exam, then the Japanese for to pass (an exam), then say \emph{gōkaku suru} again
\end{itemize}

\begin{tcolorbox}[breakable,colback=teal!4,colframe=teal!35!black,title={Before you move on}]
What does \ja{ごうかくする} mean? (\textquotedblleft{}To pass (a test); to be accepted\textquotedblright{}.) Say it once more.
\end{tcolorbox}
\section[\texorpdfstring{\ja{せいせき} (seiseki)}{せいせき (seiseki)}]{\ja{せいせき} (seiseki) — grades; results}

\label{lesson:JA-C515-seiseki}



\begin{quote}
Before the new one: say the Japanese for to pass (an exam), then the Japanese for to pass (a test).
\end{quote}

\subsection*{You'll want to know: \ja{せいせき}}
\textbf{\ja{せいせき}} — \emph{seiseki} — \textquotedblleft{}grades; results\textquotedblright{}.

\begin{quote}
\textbf{\ja{ごうかくは} \ja{せいせきに} \ja{よります。}} — \emph{gōkaku wa seiseki ni yorimasu} — \textquotedblleft{}Passing depends on your grades.\textquotedblright{}
\end{quote}

\begin{cousinweb}
\textbf{\ja{せいせき}} (\emph{seiseki}) is Sino-Japanese (kango): two Chinese characters read with their on'yomi, readings borrowed from Chinese. They are \emph{sei}, \textquotedblleft{}turn into\textquotedblright{} (Middle Chinese \emph{dzyeng}) and \emph{seki}, \textquotedblleft{}exploits\textquotedblright{} (\emph{tsek}) (KANJIDIC2; Middle Chinese in Baxter's transcription).
\end{cousinweb}

\begin{grammarlens}[title={What you've built}]
One more word for saying what happens if.
\end{grammarlens}

\subsection*{Your turn}
\begin{itemize}
\raggedright
  \item \emph{Say it:} \emph{seiseki}
  \item \emph{Say it:} \emph{seiseki}, once more
  \item \emph{Say it:} \emph{gōkaku wa seiseki ni yorimasu} — Passing depends on your grades.
  \item \emph{Recall:} say the Japanese for to pass (an exam), then the Japanese for to pass (a test), then say \emph{seiseki} again
\end{itemize}

\begin{tcolorbox}[breakable,colback=teal!4,colframe=teal!35!black,title={Before you move on}]
What does \ja{せいせき} mean? (\textquotedblleft{}Grades; results\textquotedblright{}.) Say it once more.
\end{tcolorbox}
\section[\texorpdfstring{\ja{よびこう} (yobikō)}{よびこう (yobikō)}]{\ja{よびこう} (yobikō) — a cram school (for entrance exams)}

\label{lesson:JA-C515-yobikou}



\begin{quote}
Before the new one: say the Japanese for to pass (a test), then the Japanese for grades.
\end{quote}

\subsection*{You'll want to know: \ja{よびこう}}
\textbf{\ja{よびこう}} — \emph{yobikō} — \textquotedblleft{}a cram school (for entrance exams)\textquotedblright{}.

\begin{quote}
\textbf{\ja{もし} \ja{しけんに} \ja{おちたら、よびこうに} \ja{かよいます。}} — \emph{moshi shiken ni ochitara, yobikō ni kayoimasu} — \textquotedblleft{}If I fail the exam, I will go to a cram school.\textquotedblright{}
\end{quote}

\begin{cousinweb}
\textbf{\ja{よびこう}} (\emph{yobikō}) is Sino-Japanese (kango): three Chinese characters read with their on'yomi, readings borrowed from Chinese. They are \emph{yo}, \textquotedblleft{}beforehand\textquotedblright{} (Middle Chinese \emph{yoH}), \emph{bi}, \textquotedblleft{}equip\textquotedblright{} (\emph{bijH}) and \emph{kō}, \textquotedblleft{}school\textquotedblright{} (KANJIDIC2; Middle Chinese in Baxter's transcription).
\end{cousinweb}

\begin{grammarlens}[title={What you've built}]
One more word for saying what happens if.
\end{grammarlens}

\subsection*{Your turn}
\begin{itemize}
\raggedright
  \item \emph{Say it:} \emph{yobikō}
  \item \emph{Say it:} \emph{yobikō}, once more
  \item \emph{Say it:} \emph{moshi shiken ni ochitara, yobikō ni kayoimasu} — If I fail the exam, I will go to a cram school.
  \item \emph{Recall:} say the Japanese for to pass (a test), then the Japanese for grades, then say \emph{yobikō} again
\end{itemize}

\begin{tcolorbox}[breakable,colback=teal!4,colframe=teal!35!black,title={Before you move on}]
What does \ja{よびこう} mean? (\textquotedblleft{}A cram school (for entrance exams)\textquotedblright{}.) Say it once more.
\end{tcolorbox}
\section[(dialogue)]{Looking back: from By Chance to If I Pass}

\label{lesson:JA-R515-review-510-515}



\begin{quote}
Say the Japanese for by chance and a chance.
\end{quote}

\subsection*{Your turn}
Say these aloud:

\begin{itemize}
\raggedright
  \item by chance, a chance, out of the blue, to make good use of, to be blessed (with)
  \item to cause, an achievement, to arise, how well something works, the origin
  \item to change, to switch (over), to be brought forward, to clash (of plans), changing every day
  \item to plan (an event), to look over again, to bring forward (a date), to adjust, arrangements
  \item to guess, to predict, a prospect, a forecast, an outlook
  \item to sit an entrance exam, to pass (an exam), to pass (a test), grades, a cram school (for entrance exams)
\end{itemize}

\begin{tcolorbox}[breakable,colback=teal!4,colframe=teal!35!black,title={Before you move on}]
By chance? (\textbf{\ja{ぐうぜん}}, \emph{gūzen}.) A chance? (\textbf{\ja{きっかけ}}, \emph{kikkake}.) Out of the blue? (\textbf{\ja{ふと}}, \emph{futo}.) To make good use of? (\textbf{\ja{いかす}}, \emph{ikasu}.) To be blessed (with)? (\textbf{\ja{めぐまれる}}, \emph{megumareru}.) To cause? (\textbf{\ja{ひきおこす}}, \emph{hikiokosu}.) An achievement? (\textbf{\ja{せいか}}, \emph{seika}.) To arise? (\textbf{\ja{しょうじる}}, \emph{shōjiru}.) How well something works? (\textbf{\ja{ききめ}}, \emph{kikime}.) The origin? (\textbf{\ja{もと}}, \emph{moto}.) To change? (\textbf{\ja{へんかする}}, \emph{henka suru}.) To switch (over)? (\textbf{\ja{きりかえる}}, \emph{kirikaeru}.) To be brought forward? (\textbf{\ja{はやまる}}, \emph{hayamaru}.) To clash (of plans)? (\textbf{\ja{かさなる}}, \emph{kasanaru}.) Changing every day? (\textbf{\ja{ひがわり}}, \emph{higawari}.) To plan (an event)? (\textbf{\ja{きかくする}}, \emph{kikaku suru}.) To look over again? (\textbf{\ja{みなおす}}, \emph{minaosu}.) To bring forward (a date)? (\textbf{\ja{くりあげる}}, \emph{kuriageru}.) To adjust? (\textbf{\ja{ちょうせいする}}, \emph{chōsei suru}.) Arrangements? (\textbf{\ja{だんどり}}, \emph{dandori}.) To guess? (\textbf{\ja{すいそくする}}, \emph{suisoku suru}.) To predict? (\textbf{\ja{よそくする}}, \emph{yosoku suru}.) A prospect? (\textbf{\ja{みこみ}}, \emph{mikomi}.) A forecast? (\textbf{\ja{よほう}}, \emph{yohō}.) An outlook? (\textbf{\ja{みとおし}}, \emph{mitōshi}.) To sit an entrance exam? (\textbf{\ja{じゅけんする}}, \emph{juken suru}.) To pass (an exam)? (\textbf{\ja{うかる}}, \emph{ukaru}.) To pass (a test)? (\textbf{\ja{ごうかくする}}, \emph{gōkaku suru}.) Grades? (\textbf{\ja{せいせき}}, \emph{seiseki}.) A cram school (for entrance exams)? (\textbf{\ja{よびこう}}, \emph{yobikō}.)
\end{tcolorbox}
